\section{Referencias}

\begin{itemize}
    \item Biggs, J., \& Tang, C. (2011). \textit{Teaching for quality learning at university}. McGraw-Hill Education.
    \item Deterding, S., Dixon, D., Khaled, R., \& Nacke, L. (2011). From game design elements to gamefulness: Defining gamification. \textit{Proceedings of the 15th International Academic MindTrek Conference}.
    \item Hattie, J., \& Timperley, H. (2007). The power of feedback. \textit{Review of Educational Research}.
    \item Irvine, C. (2005). CyberCIEGE: A game for computer network security education. \textit{IEEE Security \& Privacy Magazine}.
    \item ISO/IEC. (2022). \textit{ISO/IEC 27001:2022 Information security, cybersecurity and privacy protection}.
    \item Kapp, K. M. (2012). \textit{The gamification of learning and instruction}. Pfeiffer.
    \item Kolb, D. A. (1984). \textit{Experiential learning}. Prentice Hall.
    \item Michael, D., \& Chen, S. (2006). \textit{Serious games: Games that educate, train, and inform}. Thomson Course Technology.
    \item Ng, C. Y., \& Hasan, M. K. (2024). Cybersecurity serious games development: A systematic review. \textit{Computers \& Security}.
    \item NIST. (2024). \textit{The NIST Cybersecurity Framework (CSF) 2.0}. National Institute of Standards and Technology.
    \item Prince, M. (2004). Does active learning work? A review of the research. \textit{Journal of Engineering Education}.
    \item Stallings, W., \& Brown, L. (2018). \textit{Computer security: Principles and practice}. Pearson.
    \item Verizon. (2024). \textit{Data Breach Investigations Report}.
\end{itemize}
